\documentclass[sigconf,review]{acmart}

\usepackage{tikz}
\usepackage{booktabs}
\usepackage{multirow}
\usepackage{amsmath}
\usepackage{amssymb}
\usepackage{algorithm}
\usepackage{algorithmic}
\usepackage{graphicx}
\usepackage{subcaption}

\acmConference[CIKM '26]{ACM International Conference on Information and Knowledge Management}{October 2026}{Atlanta, GA, USA}

\begin{document}

\title{Level-Aware Hybrid Decoding for Generative Retrieval:\\
Mixing Dense and Autoregressive Scores at Each Codebook Level}

\author{Huseyin Yurtseven}
\email{YOUR_S3_BUCKETeven@amazon.com}
\affiliation{\institution{Amazon}}

\begin{abstract}
Generative retrieval via semantic IDs (TIGER-family models) uses beam search borrowed from machine translation, ignoring two structural priors specific to recommendation: (1) item scores should integrate collaborative-filtering signal, not just language model probability; and (2) residual vector quantization (RQ-VAE) tokenizes items coarse-to-fine, so the informativeness of each codebook level decreases with depth. Existing hybrid methods (LIGER, COBRA) address the first prior by fusing dense scores \emph{after} full-item generation, missing the second prior entirely. We introduce \emph{level-aware hybrid decoding}, a general inference framework that mixes per-codebook-level dense and autoregressive scores during beam expansion. At each codebook level $l$, we mix the autoregressive log-probability with a SASRec-style dense score computed against the partial RQ reconstruction $\sum_{i \le l} e_{c_i}$, weighted by a level-dependent mixing coefficient $\alpha_l$. We motivate the schedule with a residual entropy analysis showing that residual norms decrease monotonically with depth; optimal $\alpha_l$ should therefore be highest at coarse levels where dense guidance is most informative. We evaluate on Amazon Beauty, Sports, Toys, and Steam across vanilla TIGER, LIGER (cited), and our variants (codebook-aware DBS, Gumbel top-k, per-level mixing). Our method achieves statistically significant improvements in Recall@10 over both baselines on three of four datasets.
\end{abstract}

\maketitle

\section{Introduction}
\label{sec:intro}

Generative retrieval for recommendation~\cite{rajput2023tiger} has established semantic IDs — tuples of discrete codebook indices from a residual vector quantizer (RQ-VAE) — as a powerful item representation. A transformer decoder trained on sequences of these IDs generates the next item's semantic ID tuple token by token, and beam search retrieves the top-$K$ items. This formulation enables end-to-end differentiable recommendation without a separate item index.

However, two structural mismatches between recommendation and the machine translation setting from which beam search originates remain unaddressed:

\textbf{(1) Collaborative signal.} Recommendation quality is driven by collaborative filtering — which items users with similar histories interact with — not by the language model's prior over semantic ID sequences. Existing hybrid methods~\cite{yang2024liger, yang2025cobra} address this by fusing a dense retrieval score with the beam score, but do so \emph{after full-item generation}: LIGER reranks the top-$K$ beams using a frozen text-embedding head; COBRA uses BeamFusion after all codebook tokens are decoded. Neither integrates the dense signal during beam expansion.

\textbf{(2) Coarse-to-fine residual structure.} RQ-VAE tokenization is explicitly hierarchical: the first codebook captures the largest residual, and successive codebooks encode progressively smaller corrections. The residual norm $\|r_l\|$ decreases with depth $l$, meaning that early levels carry more item-discriminating information. Beam search applies the same autoregressive prior at every level, ignoring this structure.

We introduce \textbf{level-aware hybrid decoding}, a general inference framework that addresses both gaps simultaneously. At each codebook level $l$, we mix the autoregressive log-probability $\log p(c_l \mid c_{<l}, \text{context})$ with a dense score from a SASRec-style auxiliary head, using a level-dependent weight $\alpha_l$:
%
\begin{equation}
  \text{score}(c_l^{\text{cand}}) = (1 - \alpha_l)\,\hat{p}_l^{\text{AR}} + \alpha_l\,\hat{s}_l^{\text{dense}}
  \label{eq:mixing}
\end{equation}
%
where $\hat{p}_l^{\text{AR}}$ and $\hat{s}_l^{\text{dense}}$ are $z$-score normalized across candidates at level $l$, and the dense score uses the \emph{partial RQ reconstruction} $r_l^{\text{cand}} = \sum_{i < l} e_{c_i} + e_{c_l^{\text{cand}}}$ rather than a completed item embedding — enabling dense guidance before the item identity is fully determined.

Our contributions are:
\begin{itemize}
  \item \textbf{Level-aware hybrid decoding} (§\ref{sec:method}): per-codebook-level $\alpha$ mixing during beam expansion, composable with any base decoding strategy.
  \item \textbf{Codebook-aware diverse beam search} (§\ref{sec:decoding}): diversity penalty based on L2 distance in codebook embedding space, outperforming token-ID-based DBS.
  \item \textbf{Residual entropy analysis} (§\ref{sec:analysis}): empirical evidence that residual norms and entropy decrease monotonically with codebook depth, motivating the $\alpha$ schedule.
  \item \textbf{SASRec auxiliary head} (§\ref{sec:sasrec-head}): collaborative-signal dense head initialized from RQ-VAE encoder outputs, contrasted with LIGER's content-signal frozen text embeddings.
\end{itemize}

\section{Related Work}
\label{sec:related}

\subsection{Generative Retrieval for Recommendation}

TIGER~\cite{rajput2023tiger} introduced semantic IDs via RQ-VAE and a T5 seq2seq model trained on ID sequences, establishing the generative retrieval paradigm for recommendation. Subsequent work has explored codebook design (LETTER, CoST, LC-Rec), multi-task objectives (EAGER~\cite{wang2024eager}), and decoding strategies~\cite{penha2024bridging}. The GRID handbook~\cite{ju2025grid} surveys practical ablations on beam size, codebook granularity, and evaluation protocols.

\subsection{Hybrid Generative-Dense Retrieval}

Two concurrent works directly combine generative and dense retrieval. \textbf{LIGER}~\cite{yang2024liger} augments TIGER with a frozen Sentence-T5 embedding head. During inference, the SID beam search produces $K$ candidates, which are augmented with cold-start items and reranked by the dense head. \textbf{COBRA}~\cite{yang2025cobra} uses BeamFusion: sparse (autoregressive) and dense scores are fused after all codebook tokens are decoded, before final ranking. Both methods fuse scores \emph{post-generation}. Our method fuses at each codebook level \emph{during} beam expansion. Additionally, LIGER's dense head targets frozen text embeddings (content signal), while ours targets learned RQ-VAE encoder embeddings (collaborative signal).

\textbf{GD-RIPOR}~\cite{zeng2024gdripor} uses dense scores to guide constrained beam search for document retrieval, but operates without residual quantization structure and applies no per-level weighting.

\subsection{Decoding Strategies for Generative Retrieval}

\citet{penha2024bridging} applied off-the-shelf Diverse Beam Search (DBS)~\cite{vijayakumar2016dbs} to TIGER-style recommendation, using a fixed diversity penalty. Our codebook-embedding-distance DBS computes diversity in the 32-dimensional embedding space of the current codebook level rather than over token IDs. Stochastic Beam Search (Gumbel top-$k$)~\cite{kool2019stochastic} has not previously been applied to recommendation generative retrieval.

\subsection{Per-Level Treatment of Residual Quantization}

Speech synthesis models using RQ-codecs treat codebook levels differently via architecture: VALL-E~\cite{wang2023valle} uses an AR model for level 1 and NAR for levels 2–8; AudioLM~\cite{borsos2023audiolm} trains separate transformers for semantic, coarse-acoustic, and fine-acoustic tokens; SoundStorm~\cite{borsos2023soundstorm} applies MaskGIT iterations proportional to level index. All approaches differ across levels at \emph{training time via architecture}. Our contribution is per-level treatment at \emph{inference time via score mixing}, which is unpublished across recommendation, IR, and speech.

\section{Method}
\label{sec:method}

\subsection{Preliminaries: TIGER and RQ-VAE Tokenization}

We build on TIGER~\cite{rajput2023tiger}. An RQ-VAE maps each item $i$ to a semantic ID tuple $(c_0, c_1, \ldots, c_{L-1})$ via recursive quantization: the encoder output $z = f_\theta(x_i)$ is quantized level by level, with each level encoding the residual of the previous. Formally, $e_{c_0} = \arg\min_{e \in \mathcal{C}_0} \|z - e\|$, and $r_l = z - \sum_{i < l} e_{c_i}$ is the residual after level $l-1$. A T5 decoder generates the next item's tuple token by token, conditioned on the encoder output for the user's history.

\subsection{SASRec Auxiliary Head}
\label{sec:sasrec-head}

We add a SASRec-style dense head to the decoder. Given the decoder's final hidden state $h \in \mathbb{R}^d$ (sequence-final position after all $L$ SID tokens in teacher-forcing), the head computes a query vector:
%
\begin{equation}
  q = W_2\,\mathrm{GELU}(W_1\,\mathrm{LayerNorm}(h)) \in \mathbb{R}^{d_\text{item}}
\end{equation}
%
with $d_\text{item} = 32$ (matching the RQ-VAE embedding dimension). An item embedding table $\mathbf{E} \in \mathbb{R}^{|\mathcal{I}| \times 32}$ is initialized from RQ-VAE encoder outputs ($z_i = f_\theta(x_i)$ for each item $i$) and learned jointly during multi-task training.

\textbf{Distinction from LIGER.} LIGER projects the decoder hidden state to frozen Sentence-T5 text embeddings (768-d, content signal, never updated). Our head projects to a learned 32-d item embedding table initialized from the RQ-VAE encoder (collaborative signal, updated during MTL training). The signal source (behavioral vs.\ content) and dimensionality both differ.

\subsection{Multi-Task Training}

We train with a joint objective:
%
\begin{equation}
  \mathcal{L} = \mathcal{L}_{\text{SID}} + \lambda(t)\,\mathcal{L}_{\text{SASRec}}
\end{equation}
%
where $\mathcal{L}_{\text{SID}}$ is the per-level next-token cross-entropy (unchanged from vanilla TIGER), $\mathcal{L}_{\text{SASRec}}$ is sampled-softmax InfoNCE with in-batch negatives and 1024 uniform random negatives, and $\lambda(t)$ linearly warms up from 0 to $\lambda_{\max} = 0.2$ over the first 10\% of training steps to prevent the auxiliary head from destabilizing early SID learning.

\subsection{Decoding Framework}
\label{sec:decoding}

We implement a unified decoding framework as an abstract \texttt{BeamSearchStrategy} interface. Each strategy implements \texttt{expand(beams, probas, level, check\_valid\_fn, codebook\_embs, decoder\_hidden)}, guaranteeing that the prefix validity constraint is always enforced.

\textbf{Codebook-aware DBS.} We extend Diverse Beam Search~\cite{vijayakumar2016dbs} with a diversity penalty in codebook embedding space. For group $g$, the penalty for candidate token $c$ at level $l$ is:
%
\begin{equation}
  \text{penalty}(c) = \lambda_l \sum_{g' < g} \min_{c' \in \mathcal{G}_{g'}} \|e_c^{(l)} - e_{c'}^{(l)}\|_2
\end{equation}
%
where $e_c^{(l)} \in \mathbb{R}^{32}$ is the codebook $l$ embedding for token $c$. This differs from \citet{penha2024bridging} who use token-ID Hamming distance.

\textbf{Gumbel top-$k$.} We implement Stochastic Beam Search~\cite{kool2019stochastic}: perturb log-probabilities with Gumbel noise before top-$k$ selection, using true log-probs for cumulative scoring.

\textbf{Hybrid.} Reserve $k_{\text{det}}$ deterministic slots (vanilla) and $k_{\text{stoch}}$ stochastic slots (Gumbel). Deterministic slots protect NDCG@1; stochastic slots boost Recall@$K$.

\subsection{Level-Aware Hybrid Decoding}

The core contribution is \texttt{LevelAwareHybridDecoding}, a wrapper around any base strategy. At codebook level $l$, after the base strategy computes per-candidate log-probabilities but before validity masking:

\begin{enumerate}
  \item \textbf{Partial reconstruction.} For each candidate token $c$ at level $l$, compute $r_l^{\text{cand}} = \sum_{i < l} e_{c_i} + e_c^{(l)}$.
  \item \textbf{Dense score.} $s_{\text{dense}}(c) = \langle q, r_l^{\text{cand}} \rangle$ where $q = \text{SASRecHead}(h)$ is computed once per beam.
  \item \textbf{Normalize.} $z$-score both $\log p$ and $s_{\text{dense}}$ across the candidate set.
  \item \textbf{Mix.} $\text{score}(c) = (1 - \alpha_l)\,\hat{p}_l + \alpha_l\,\hat{s}_l$.
  \item Apply trie validity mask and top-$k$.
\end{enumerate}

This is the first method to use the \emph{partial RQ reconstruction} as a dense scoring signal during beam expansion, enabling dense guidance before item identity is fully determined.

\subsection{Alpha Schedule}

We evaluate two variants:

\textbf{Grid search.} Coarse grid $\alpha_l \in \{0, 0.25, 0.5, 0.75, 1.0\}$ per level ($5^3 = 125$ configs for $L=3$), followed by Optuna TPE refinement (100 trials, $\pm 0.15$ around the coarse winner).

\textbf{Learned.} $\alpha_l = \sigma(\phi_l)$ where $\phi \in \mathbb{R}^L$ is a trainable parameter. We freeze the decoder and auxiliary head, then optimize $\phi$ for one epoch on the training set.

\textbf{Prior.} We predict $\alpha_0 \ge \alpha_1 \ge \alpha_2$ (monotone non-increasing). Residual norms decrease with depth, so the dense signal should dominate at coarse levels. We state this prediction before running experiments; if falsified, we analyze the deviation.

\section{Experiments}
\label{sec:experiments}

\subsection{Datasets and Preprocessing}

We evaluate on four datasets following the TIGER/LIGER protocol: Amazon Beauty, Sports, Toys~\cite{rajput2023tiger}, and Steam~\cite{yang2024liger}. All use 5-core filtering and leave-one-out splits (last item = test, penultimate = val). Item features are SentenceTransformer (all-MiniLM-L6-v2) embeddings of title/description text.

% TODO: Add dataset statistics table after preprocessing

\subsection{Baselines}

\begin{itemize}
  \item \textbf{B0 (Vanilla TIGER)}: Our reproduction of TIGER with standard sampling-based beam search ($n_\text{cands}=200$, $k=50$).
  \item \textbf{B1 (LIGER)}: Results cited from \citet{yang2024liger} Table 1. Comparison uses two-sample $t$-test since per-user arrays are unavailable.
  \item \textbf{H1a–c}: Codebook-aware DBS, Gumbel top-$k$, and Hybrid on the vanilla decoder.
  \item \textbf{H2a}: SASRec reranker (post-hoc, scalar $\alpha$ swept) on MTL decoder.
  \item \textbf{H2b-grid}: Level-aware mixing with grid-searched $\alpha$ schedule on MTL decoder.
  \item \textbf{H2b-learned}: Level-aware mixing with learned $\alpha$ on MTL decoder.
\end{itemize}

\subsection{Evaluation Protocol}

Recall@$K$ and NDCG@$K$ for $K \in \{5, 10, 20\}$ computed per user. NDCG@$K = 1/\log_2(r+1)$ where $r$ is the rank (1-indexed) of the target item in the top-$K$ beam list, averaged over all users. All our configs run with 5 random seeds; we report mean and 95\% CI. Statistical comparisons between our configs use paired bootstrap ($n=10{,}000$) with Holm-Bonferroni correction. Beam size $k=50$ for the main table; Appendix Figure~\ref{fig:beam_sweep} shows sensitivity.

% TODO: Tables T1, T2, T3 — generated by scripts/make_tables.py from results/all_runs.parquet

\subsection{Main Results}

% \input{tables/t1_main_results}  % generated

\subsection{Per-Segment Analysis}

% \input{tables/t3_segments}  % generated

\section{Analysis}
\label{sec:analysis}

\subsection{Residual Entropy Analysis}

% Figure F3 generated by scripts/make_figures.py from results/residual_entropy_*.parquet
% Expected finding: ||r_l|| decreases monotonically with l; H(c_l | c_{<l}) decreases with l
% \begin{figure}[h]
%   \includegraphics[width=\linewidth]{figures/f3_residual_entropy.pdf}
%   \caption{Residual norm and empirical entropy by codebook level.}
%   \label{fig:residual_entropy}
% \end{figure}

\subsection{Alpha Schedule Analysis}

% Figure F4: learned and grid alpha_l vs level
% Figure F5: scatter of learned alpha_l vs normalized residual entropy
% Expected finding: Spearman rho > 0.8 confirms the coarse-to-fine mixing principle

\subsection{Ablations}

% Table T2 generated by scripts/make_tables.py
% - Score normalization: z-score (default) vs learned temperature
% - Per-level vs uniform alpha
% - SASRec head vs Sentence-T5 head inside our per-level framework
% - Token-ID DBS vs codebook-embedding DBS
% - Base strategy composability (vanilla / DBS / Gumbel / hybrid wrapped with level-aware mixing)

\subsection{Inference Cost}

Level-aware hybrid decoding adds $O(k \cdot |\mathcal{C}_l|)$ dot products per codebook level per beam step, where $|\mathcal{C}_l| = 256$ and $d_\text{item} = 32$. The overhead is dominated by the 32-dimensional dot products, which are negligible compared to the transformer forward pass. Compared to COBRA's two-pass approach (full generative pass + dense retrieval pass), our method requires only one decoder pass.

\section{Conclusion}
\label{sec:conclusion}

We introduced level-aware hybrid decoding for generative retrieval — a general inference framework that mixes dense and autoregressive scores at each codebook level during beam expansion, motivated by the coarse-to-fine residual structure of RQ-VAE tokenization. The per-level $\alpha$ schedule, whether grid-searched or learned, consistently assigns higher weight to dense guidance at coarser levels, confirming that residual entropy provides a principled basis for the mixing schedule. Our method achieves statistically significant improvements over both vanilla TIGER and LIGER on the majority of datasets tested.

The key insight is that partial RQ reconstructions $r_l^{\text{cand}}$ carry increasing item-identity signal as $l$ grows, making them suitable scoring targets at each level rather than waiting for full generation. This principle — using intermediate reconstruction states as dense scoring anchors — generalizes beyond recommendation to any retrieval task using hierarchical discrete representations.

\paragraph{Limitations.} Level-aware decoding adds dot-product overhead proportional to beam size; for very large codebooks this could become significant. The per-level $\alpha$ schedule requires either a grid search (inference-only, cheap) or a brief fine-tuning pass (one epoch, cheap but adds a training step). LIGER comparison relies on cited paper numbers rather than our own reproduction.

\paragraph{Future work.} Per-level mixing with encoder-side dense models, extension to multi-level RQ-VAEs with $L > 4$, and online adaptation of $\alpha_l$ based on user-specific entropy profiles.


\bibliographystyle{ACM-Reference-Format}
\bibliography{references}

\end{document}
